\section{The record's ledger and its central claims}\label{sec:ledger}

The record labels each claim with its own status: \textsf{VERIFIED} for a complete proof, \textsf{OPEN} when neither a proof nor a counterexample is recorded, and \textsf{REFUTED} for an exact assertion with a recorded contradiction or counterexample. These are the AI verification's labels. Two lists carry them, and this reader cites both:
\begin{itemize}
\item the \emph{claim ledger} in the evidence archive \texttt{04\_}, with 48 rows: 42 \textsf{VERIFIED}, 2 \textsf{OPEN}, 4 \textsf{REFUTED}. The README quotes its counts.
\item the monograph's \emph{status list}, the 21 numbered items of its Part~I. Item numbers below refer to this list.
\end{itemize}

\subsection{The central claims (record: \textsf{VERIFIED})}

\begin{itemize}
\item \emph{The construction.} According to the monograph, its sections 6--7 prove the period-family descent, the toric and finite fillings, the global gluing, and the integral nearby-cycle and Leray comparisons. With them the constructed threefold $X$ satisfies $H^1(X;\ZZ)=H^2(X;\ZZ)=0$, $a(X)=1$ and $c_3(X)=2$, and $f:X\to\Pp^1$ is holomorphic and an algebraic reduction.
\item \emph{A counterexample to the published theorem.} $X$ satisfies every stated hypothesis of Theorem~2.2 of \cite{CDP20} and contradicts its conclusion (monograph §13.11, item~5). The monograph says that this uses the complete construction and topology proofs, not only the two defects below.
\item \emph{A corrected criterion.} This is the part of the published argument that survives. If some line bundle $L_0$ has $H^2(X,T_X\otimes L_0)=0$, then generic vanishing holds and $c_3(X)\le0$ (monograph Theorem 13.28; the correction note's Theorem 5.2). For the constructed $X$ this hypothesis fails for every $L_0$, since $c_3(X)=2$.
\item \emph{Further invariants}: the Hodge diamond and the Frölicher sequence, the Bott--Chern and Aeppli dimensions, and global representatives with the Bott--Chern product structure (status list, items 17, 19 and 20).
\end{itemize}
\status{\textsf{Located}. Not checked here. They are the substance of the record, and they need specialist verification; the record itself lists that adjudication as open.}

\subsection{The four rows labelled \textsf{REFUTED}}

Each row is given with its claim-ledger identifier.

\begin{enumerate}[label=\arabic*.]
\item \emph{The singular-fibre inference of \cite{CDP20}} (p.~692 of the journal version; \texttt{S6-CDP-P692}; status list, item~7).
\begin{itemize}
\item The inference: for an irreducible reduced fibre, the published proof reduces a vanishing to the vanishing of $H^0(S,\widetilde\Omega^1_S\otimes L|_S)$.
\item The record says that this is not a formal consequence of the two cotangent sequences, and gives a test family.
\item The source manuscript had already located the failure at its own cusp fibre (its Lemma~10.3 and Theorem~10.5), and had noted that the step is valid on a normal fibre (its Remark~10.4).
\end{itemize}
\status{The test family and the corrected criterion are \textsf{audited} here. The exact size of the failure at fibres with normal crossings, and the exact condition for the torsion that causes it, are \textsf{proved here} (Section~\ref{sec:correction}). The source manuscript quotes the published sentence verbatim (its §10.1). The paper itself was not re-read here.}
\item \emph{A homology formula imported from \cite{CDP98}} (\texttt{S6-CDP98-MONO}; status list, item~12). The direct-product formula of the 1998 paper (its Lemma~3.2, in the record's reading) omits the coinvariants of the monodromy. This leaves a torus branch of the 2020 argument without its printed rank argument. \status{\textsf{Located}.}
\item \emph{A remark of the source manuscript on relative differentials} (\texttt{S6-HODGE-CUSP-TORSION}; status list, item~18).
\begin{itemize}
\item The manuscript's Remark~1.3 speaks of the sheaves $\Omega^p_{X/B}$, and its Remark~9.21 of $\Omega^1_{X/B}$, as having torsion along all three singular fibres.
\item The record shows that at the cusp $\Omega^1_{X/B}$ is torsion-free but not locally free, and that its torsion occurs only at the two multiple fibres.
\end{itemize}
\status{\textsf{Proved here} for the local models the record gives (Proposition~\ref{prop:torsion}). The refutation is correct for $p=1$. For $p=2$ the same local models do have torsion along the double curves of the cusp fibre, so Remark~1.3's statement holds for $\Omega^2_{X/B}$.}
\item \emph{A remark of the source manuscript on the gluing integers} (\texttt{S6-LOCAL-728}; monograph §14, item~2). The remark (its Remark~7.28) is stated for a fixed triple. The triple $(\ell_0,\ell_1,\ell_2)=(100,1,-1)$ gives $12\ell_0-4\ell_1-3\ell_2=1199\neq\pm1$, so the statement must quantify existentially over $\ell_0$ for fixed admissible $(\ell_1,\ell_2)$. \status{\textsf{Reproduced}: the arithmetic.}
\end{enumerate}

\subsection{The two rows labelled \textsf{OPEN}}

\begin{itemize}
\item \texttt{S6-OPEN-HODGE-REPLICATION}: independent specialist replication of the cyclic invariant lattices, the logarithmic determinant divisor, the residue cancellation and an order-four scalar test. The monograph's status list has this as its single open item (item~21), together with the three cotangent direct images and the values $h^{1,2}=1$, $h^{1,1}=2$.
\item \texttt{S6-OPEN-EXTERNAL-ADJUDICATION}: specialist and community adjudication of the complete construction and counterexample.
\end{itemize}
The monograph states that these rows are open because independent human replication and adjudication are outstanding, and that they are not unresolved premises of its proofs.

\subsection{Confirmed local repairs}

Monograph §14 lists seven items, labelled separately from the theorem. Five are labelled \textsf{VERIFIED}:
\begin{itemize}
\item a closed-image argument;
\item an appendix sequence whose arrows are cohomological although it is labelled as homology;
\item a fundamental-group identification;
\item a build defect of the record's own assembler;
\item a correction to a finite-filling computation of a canonical section.
\end{itemize}
The other two are the \textsf{REFUTED} remarks treated above. The five are \textsf{located}.
